\section{Method}
\label{sec:method}

%% Pipeline figure (TikZ, single column, drawn at native width ~7.9 cm).
\begin{figure}[t]
\centering
\begin{tikzpicture}[
  font=\sffamily\scriptsize,
  box/.style={draw, rounded corners=2pt, align=center, text width=1.62cm, minimum height=8mm, inner sep=2pt, fill=white},
  frozen/.style={box, fill=gray!15},
  anch/.style={box, fill=blue!8},
  prob/.style={box, fill=orange!18},
  outbox/.style={box, text width=1.75cm, fill=white, very thick},
  opt/.style={box, text width=1.75cm, dashed, fill=green!6},
  arr/.style={-{Stealth[length=1.6mm]}, semithick},
  darr/.style={-{Stealth[length=1.6mm]}, semithick, dashed}]
  % column 1: inputs
  \node[box]    (rgb)    at (0.85, 2.15) {RGB image $x$};
  \node[box]    (sensor) at (0.85, 0.45) {sparse depth,\\budget $k$};
  % column 2: model and split
  \node[frozen] (f)      at (2.85, 2.15) {frozen model $f$\\[-1pt]{\tiny(or any probe-indep.\ centre)}};
  \node[anch]   (S)      at (2.85, 0.95) {$k{-}n$ structured\\anchors $\cS$};
  \node[prob]   (Q)      at (2.85,-0.15) {$n$ uniform\\probes $\cQ$};
  % column 3: centre/scale and threshold
  \node[anch]   (ms)     at (4.85, 2.15) {centre $m$, scale $s_\gamma$\\[-1pt]{\tiny from $\cS$ only}};
  \node[prob]   (q)      at (4.85,-0.15) {threshold\\$\qhat_i=R_{(r)}$};
  % column 4: output and optional stages
  \node[outbox] (int)    at (6.95, 2.15) {$\yhat\,e^{\,m\pm\qhat_i s_\gamma}$\\[-1pt]\textbf{\mDual}};
  \node[opt]    (hull)   at (6.95, 0.95) {$\cup$ \mDistSplit\\[-1pt]\textbf{\mDualBoth}};
  \node[opt]    (env)    at (6.95,-0.15) {sensor-error\\envelope};
  % arrows
  \draw[arr] (rgb) -- (f);
  \draw[arr] (f) -- node[above, font=\tiny\sffamily]{$\yhat$} (ms);
  \draw[arr] (sensor.east) -- ++(0.12,0) |- (S.west);
  \draw[arr] (sensor.east) -- ++(0.12,0) |- (Q.west);
  \draw[arr] (S.east) -- (ms.south west);
  \draw[arr] (Q) -- (q);
  \draw[arr] (ms) -- (q);
  \draw[arr] (ms) -- (int);
  \draw[arr] (q.east) -- ++(0.12,0) |- ([yshift=-2mm]int.west);
  \draw[darr] (int) -- (hull);
  \draw[darr] (hull) -- (env);
  \node[font=\tiny\sffamily, text=orange!70!black, anchor=north] at (2.85,-0.62) {never seen by $f$, $m$, $s$};
\end{tikzpicture}
\caption{\textbf{Dual-role calibration.} Structured anchors $\cS$ build a centre $m$ and scale $s_\gamma$ around a frozen prediction $\yhat$ (or any probe-independent centre); the $n$ uniform probes $\cQ$ only score it and set one threshold per image (\cref{thm:dual}). Dashed: optional stages (\ifenvelopemain\cref{thm:both,prop:env}\else\cref{thm:both}; Supp.~D\fi).}
\label{fig:pipeline}
\end{figure}


\subsection{Architecture}
\label{sec:arch}
In the pipeline of \cref{fig:pipeline},
(1)~a frozen monocular model produces $\yhat$; nothing is retrained.
(2)~The sensor budget is split into $k-n$ anchors that follow the deployed layout and $n$ uniformly random probes.
(3)~The anchors alone build a local centre $m(p)$ and scale $s_\gamma(p)$ around $\yhat$ (\cref{sec:centre}).
(4)~The probes, unseen by stages 1--3, set one conformal threshold per image (\mDual, \cref{sec:dual}).
Optional stages add a per-distance-bin guarantee (\mDualBoth, \cref{sec:both}) and convert coverage of sensor readings into coverage of depth (\envref).
Stage~1 can be any predictor that does not see the probes, \eg a depth-completion network fed RGB and $\cS$; stage~3 is then refit around its output, so \mDual is a calibration layer, not a depth estimator.

\subsection{Centre and scale from structured anchors}
\label{sec:centre}
We use the distance-aware anchor model of the baseline \mDistSplit, fitted on the structured anchors only (details in Supp.~E).
For pixel $p$, the local shift $m(p)$ is one of four estimates from the anchor residuals $\dtil|_\cS$: the image median, the median of the 16 nearest anchors, or a Gaussian or bilateral (also weighting similarity in $\log\yhat$) kernel mean.
The local spread $\sigma(p)$ is the median absolute deviation of the 16 nearest anchor residuals, $\delta_\cS(p)$ is the distance to the nearest structured anchor as a fraction of the image diagonal, and
\begin{equation}
s(p)=\big(\sigma(p)+\tau\big)\,g\big(\delta_\cS(p)\big),
\label{eq:scale}
\end{equation}
with a floor $\tau$ and a non-decreasing $g$ fitted by isotonic regression of normalized error against distance.
The variant, bandwidths, $\tau$, $g$ and a cap on $s$ are chosen per layout on held-out \emph{head-train} images, never on calibration or test data.
The only property the guarantee below needs is that $m$ and $s$ are functions of $(x,\yhat,\cS,\dtil|_\cS)$, \ie that they do not depend on the probes.

\subsection{Per-image threshold from random probes (\mDual)}
\label{sec:dual}
When the anchors set the shift wrong everywhere, as a single cluster can, the residual is nearly flat, the threshold is set by the probes with the smallest $s$, and widths elsewhere grow by up to $\max s/\min s$.
We therefore shrink the scale toward its per-image median $M_i=\med_{p\in\cP_i}s(p)$:
\begin{equation}
s_\gamma(p)=M_i\,\big(s(p)/M_i\big)^{\gamma},\qquad \gamma=0.5 .
\label{eq:shrink}
\end{equation}
$s_\gamma$ depends on $s$ only, so it is as probe-independent as $s$.
For probes $\cQ=\{p_1,\dots,p_n\}$ with reported residuals $\dtil(p_j)$, define the scores and threshold
\begin{equation}
R_j=\frac{|\dtil(p_j)-m(p_j)|}{s_\gamma(p_j)},\quad
\qhat_i=R_{(r)},\quad r=\lceil(n+1)(1-\alpha)\rceil,
\label{eq:threshold}
\end{equation}
where $R_{(r)}$ is the $r$-th smallest score ($\qhat_i=\infty$ if $r>n$).
The interval at every evaluation pixel is
\begin{equation}
C(p)=\Big[\yhat(p)\,e^{\,m(p)-\qhat_i s_\gamma(p)},\;\yhat(p)\,e^{\,m(p)+\qhat_i s_\gamma(p)}\Big].
\label{eq:interval}
\end{equation}

\begin{theorem}[Dual-role validity]
\label{thm:dual}
Fix an image with valid pixels $\cP$ and a structured anchor set $\cS\subset\cP$ chosen by any rule, which may depend on $x$, $\yhat$ and the anchor values.
Let $m$, $s>0$ and $\gamma$ be determined by $(x,\yhat,\cS,\dtil|_\cS)$.
Let $v(p)$ be the log residual of the quantity the probes report ($v=d$ for exact depth, $v=\dtil$ for a sensor), defined on $\cU=\cP\setminus\cS$.
Draw $\cQ$ uniformly without replacement from $\cU$, compute $\qhat$ by \cref{eq:threshold} with $v$ in place of $\dtil$, and let $J$ be uniform on $\cU\setminus\cQ$.
Then
\begin{equation*}
\P\big(\yhat(J)e^{v(J)}\in C(J)\,\big|\,x,\cS,\dtil|_\cS\big)\;\ge\;\frac{r}{n+1}\;\ge\;1-\alpha,
\end{equation*}
with equality on the left when the scores have no ties and $r\le n$.
Equivalently, $\E[\Cov_i\mid x,\cS]\ge r/(n+1)$, the expectation taken over the probe draw.
\end{theorem}

\noindent\emph{Proof sketch.}
Given $(x,\cS,\dtil|_\cS)$, every $p\in\cU$ carries a fixed score $R(p)=|v(p)-m(p)|/s_\gamma(p)$, and the probes with $J$ form a uniformly random $(n+1)$-subset of $\cU$, so $R(J)$ is among the $r$ smallest of the $n+1$ scores, hence covered, with probability at least $r/(n+1)$. Supp.~A. \hfill$\square$

This is the split argument of~\cite{gui2023cmc,mao2024spatial} applied inside one image around a frozen prior; we claim no new theory.
It needs no calibration images and uses all $n$ probes for the threshold ($n\ge 9$ suffices at $\alpha=0.1$), whereas \mLocalSplit spends half on a shift.
\emph{Scope:} the guarantee holds in expectation over the probe draw, not per region; probes must be uniform over the pixels it concerns (\eg where the sensor returns); every choice, $\gamma$ included, is fixed without probe values; and it covers $v$, what the probes measure.

\subsection{Probe budget}
\label{sec:budget}
\Cref{thm:dual} fixes the mean of $\Cov_i$; $n$ fixes its spread.

\begin{proposition}[Per-image coverage law]
\label{prop:budget}
Under \cref{thm:dual} with no ties and $r\le n$, let $N=|\cU|$.
The number of covered evaluation pixels, $(N-n)\Cov_i$, follows exactly the beta-binomial law $\BetaBin(N-n,\,r,\,n+1-r)$.
As $N\to\infty$, $\Cov_i\to\BetaD(r,\,n+1-r)$ in distribution.
\end{proposition}

\noindent This is the finite-population form of the Beta law of split-conformal coverage conditional on the calibration set~\cite{vovk2012conditional}; with thousands of evaluation pixels per image the Beta approximation is excellent.
That measurements match this law (\cref{sec:res:budget}) is an implementation check implied by the randomization, not a finding.

\noindent\fbox{\parbox{\dimexpr\linewidth-2\fboxsep-2\fboxrule\relax}{\textbf{Design rule.}
The share of images below coverage $c$ in one draw is $\P(\BetaD(r,n{+}1{-}r)<c)=\P(\mathrm{Bin}(n,c)\ge r)$, known before any data are seen (\cref{tab:budget}).
At $\alpha=0.1$ it falls slowly with $n$ ($c=0.85$: \NumRuleShortHundred\% at $n=100$); the nominal level is the stronger lever: $n=r=20$ (nominal \NumRuleNomTwentyAll) leaves \NumRuleShortTwentyAll\% below 0.85 (measured on NYU and KITTI: \NumRmaxShort, coverage \NumRmaxCovLo--\NumRmaxCovHi), at \NumRmaxWidthCostLo--\NumRmaxWidthCostHi\% more width and \NumRmaxISCostLo--\NumRmaxISCostHi\% more interval score than $r=19$.}}

\begin{table}[t]
\centering\small
\setlength{\tabcolsep}{4pt}
\begin{tabular}{@{}rrccc@{}}
\toprule
$n$ & $r$ & $\frac{r}{n+1}$ & p10 of $\Cov_i$ & $\P(\Cov_i<0.85)$, \% \\
\midrule
10 & 10 & 0.909 & \NumBetaPTenTen   & \NumBetaShortOneTen \\
20 & 19 & 0.905 & \NumBetaPTenTwenty & \NumBetaShortOneTwenty \\
40 & 37 & 0.902 & \NumBetaPTenForty & \NumBetaShortOneForty \\
\bottomrule
\end{tabular}
\caption{\textbf{Probe budget at $\alpha=0.1$, one probe draw} (Beta approximation of \cref{prop:budget}; from the definition alone, no data).}
\label{tab:budget}
\end{table}

\subsection{Adding a per-distance-bin guarantee (\mDualBoth)}
\label{sec:both}
\mDual's coverage, in practice, decreases with distance to the nearest anchor.
The baseline \mDistSplit targets this: it fits the model of \cref{sec:centre} on all $k$ measurements and calibrates one threshold per bin of $\delta$ by conformal risk control~\cite{angelopoulos2024crc} over calibration \emph{scenes}, the loss being the miscovered fraction of bin pixels (a Mondrian construction~\cite{vovk2005alrw}; sparse bins are merged, Supp.~C).
For exchangeable scenes and the same anchor rule at calibration and test, the expected bin miscoverage of a new scene is at most $\alpha$ in every (merged) bin.
Unlike \mDual, this needs calibration scenes with dense depth from the same source as the measurements, and it certifies that source.
\mDualBoth is the pointwise hull
$C^{\text{both}}(p)=\big[\min(\ell^{\text{dual}},\ell^{\text{dist}}),\,\max(u^{\text{dual}},u^{\text{dist}})\big]$.

\begin{theorem}[Hull validity]
\label{thm:both}
\mDualBoth has the per-image guarantee of \cref{thm:dual} and the per-bin guarantee of \mDistSplit.
\end{theorem}

\noindent\emph{Proof sketch.} The hull contains each interval pointwise and both statements are monotone in the interval~\cite{gasparin2024merging}. Supp.~C. \hfill$\square$

\mDualBoth replaces \mDistBoth, the hull of \mDistSplit with \mLocalSplit. Any hull costs width, most when the two centres disagree, as under clustered anchors.

\ifenvelopemain
\subsection{Optional: sensor-error envelope}
\label{sec:env}
\TODO{gate: iPad-probe truth coverage $\ge0.9$ and IS $\ge$10\% below \mDistBoth, else \texttt{\textbackslash envelopemainfalse}.}
%% Sensor-error envelope: shared body of main Sec. 4.6 (\envelopemaintrue) and Supp. D (\envelopemainfalse).
%% Compiled in both documents, so it refers to Theorem 1 by its fixed main-paper number, not by \cref.
With sensor-valued probes, Theorem~1 certifies the reading $\ytil$, not the depth $y$.
Let $e(p)=\log\ytil(p)-\log y(p)$ and suppose a reference (here a laser scan) exists on calibration scenes.
We (a) subtract a global bias $\hat\beta$, the median of $e$ on head-train pixels, from the probe values, and (b) give each sensor-confidence class $c$ (ARKit: $0/1/2$) the band $[-a_c,b_c]=[Q_c(\lambda/2),Q_c(1-\lambda/2)]$ of the corrected head-train error, clipped to contain $0$.
One conformal risk control~\cite{angelopoulos2024crc} over the $U$ calibration scenes sets $\hat\lambda$: the largest grid value with $\frac{U}{U+1}\hat L(\lambda)+\frac{1}{U+1}\le\alpha_e$, $\hat L$ the mean fraction of a scene's pixels outside their band (none feasible: infinite bands).
With $[\ell,u]$ the \mDual interval for the corrected reading at level $1-\alpha_s$, the enveloped interval is
\begin{equation}
C^{\text{env}}(p)=\big[\ell(p)\,e^{-b_{c(p)}},\;u(p)\,e^{a_{c(p)}}\big].
\label{eq:envelope}
\end{equation}

\begin{proposition}[Envelope]
\label{prop:env}
If the sensor reading and its confidence are defined at every evaluation pixel, the miscoverage of $y$ by $C^{\text{env}}$, averaged over test scenes and probe draws, is at most $\alpha_s+\alpha_e$.
\end{proposition}

\ifenvelopemain
\noindent\emph{Proof sketch.} $y=\ytil e^{-\hat\beta}e^{-(e-\hat\beta)}$, so $y\notin C^{\text{env}}(p)$ requires either the corrected reading outside $[\ell,u]$ (probability at most $\alpha_s$ per image, Theorem~1) or the corrected error outside its band (at most $\alpha_e$ in expectation, by risk control); a union bound concludes. Supp.~D. \hfill$\square$

\fi
The per-image guarantee on $y$ is not retained: the envelope's part holds on average over scenes.
We fix $\alpha_s=\alpha_e=0.05$ before the run, so that $\alpha_s+\alpha_e=\alpha$.

\else
\paragraph{Sensor-valued probes.} With sensor-valued probes, \cref{thm:dual} certifies the reading; an optional sensor-error envelope that converts this into an average-over-scenes guarantee on depth is given in Supp.~D; it restored truth coverage here (\NumEnvTruthCov{} with iPad probes) but at a higher interval score than \mDistBoth.
\fi
